% GENERATED FILE. Edit canonical lessons, then run npm run generate:books.
% canonical-source-hash: fnv1a64:90aa1eea0a9ff01d
% canonical-lessons: TE-C13-sharira-bhagalu, TE-S117-letter-va

\chapter{Head and Hand}
\label{ch:te-head-and-hand}
\hlchaptermodality{13}

\begin{chapteropening}
\textbf{By the end of this chapter:} \emph{I can name the head and the hand in Telugu.}

Say \te{తల}, plainly matching its cousins, and \te{చెయ్యి} — the same Dravidian root as kai, reshaped by sound change past easy recognition.
\end{chapteropening}

\section[tala ceyyi]{\te{తల}, \te{చెయ్యి} — head and hand}

\label{lesson:TE-C13-sharira-bhagalu}



\begin{quote}
Head matches its cousins plainly. Hand looks, at first glance, like Telugu went its own way — but look closer, and it's the same word after all, just heavily reshaped.
\end{quote}

\subsection*{You'll want to know: The words}
\begin{itemize}
\raggedright
  \item \textbf{\te{తల}} (\emph{tala}) = \textquotedblleft{}\textbf{head}\textquotedblright{} — native Dravidian, matching Tamil \emph{talai}, Malayalam \emph{thala}, and Kannada \emph{tale} closely and plainly.
  \item \textbf{\te{చెయ్యి}} (\emph{ceyyi}) = \textquotedblleft{}\textbf{hand}\textquotedblright{} — also native Dravidian, and \textbf{the very same root} as Tamil/Malayalam/Kannada's \emph{kai}: Proto-Dravidian \emph{\textbf{kay-}}. Telugu's own regular sound change turned that ancient \textbf{k-} into \textbf{c-}, and reshaped the rest of the word alongside it — so \emph{ceyyi} only \emph{looks} unrelated to \emph{kai}. It's a cognate in disguise, not a different word.
\end{itemize}

So all four Dravidian languages actually \textbf{agree} on both \textquotedblleft{}head\textquotedblright{} and \textquotedblleft{}hand\textquotedblright{} at the root level — Telugu's hand-word has simply drifted further in sound than its cousins have, thanks to a sound change that reshaped many Telugu words the same way.

\subsection*{Your turn}
Say these aloud:

\begin{itemize}
\raggedright
  \item \textquotedblleft{}tala\textquotedblright{} — head, matching its cousins plainly
  \item \textquotedblleft{}ceyyi\textquotedblright{} — hand, the SAME root as kai, reshaped by sound change
  \item the lesson — looks different, is actually the same word
\end{itemize}

\begin{tcolorbox}[breakable,colback=teal!4,colframe=teal!35!black,title={Before you move on}]
Does Telugu's word for \textquotedblleft{}head\textquotedblright{} match its cousins? (\textbf{Yes} — \emph{tala}, like Tamil/Malayalam/Kannada.) Is its word for \textquotedblleft{}hand,\textquotedblright{} \emph{ceyyi}, truly a different word from \emph{kai}? (\textbf{No} — it's the \textbf{same Proto-Dravidian root} \emph{kay-}, reshaped by Telugu's own regular k$\to$c sound change — a cognate in disguise, not an unrelated word.)
\end{tcolorbox}
\section[va]{\te{వ} — one character, met inside words you already say}

\label{lesson:TE-S117-letter-va}



\begin{quote}
Before the new one: \te{ద} — what does it do? And one from further back: \te{ఇ}?

One character this time — and you have been saying it for pages without knowing which mark on the page it was.
\end{quote}

\subsection*{Script you'll notice: \te{వ}}
\textbf{\te{వ}} — \emph{va}.

It is a \textbf{consonant}, and in this script a consonant is never bare: it comes with an \emph{a} already in it. So it is not \emph{v}, it is \textbf{va}.

You already say these, and every one of them has \te{వ} somewhere inside it:

\begin{itemize}
\raggedright
  \item \textbf{\te{ధన్యవాదములు}} \emph{dhanyavādamulu} — thank you (dhanyavādamulu — \textquotedblleft{}utterances of 'worthy'\textquotedblright{})
  \item \textbf{\te{అవును}} \emph{avunu} — yes (avunu)
  \item \textbf{\te{నువ్వు} / \te{మీరు}} \emph{nuvvu / mīru} — you (familiar / respectful)
  \item \textbf{\te{పరవాలేదు}} \emph{paravālēdu} — it's okay / no problem / you're welcome
\end{itemize}

\subsection*{Writing: \te{వ} — copy what you see}
Put your pen on \te{వ} and follow its line. Copy the shape you can see — slowly, and larger than it is printed.

\begin{quote}
This book does not yet tell you \textbf{where to start the character or which way to travel}. That is a real thing, taught with real variation from school to school, and it is not written down here until it can be written down with a source. Copying what is in front of you needs no such source, and it is how the shape gets into your hand in the meantime.
\end{quote}

\subsection*{Your turn}
\begin{itemize}
\raggedright
  \item \emph{Look:} at these words, and find \te{వ} in the ones that have it
\end{itemize}

\begin{quote}
\te{ధన్యవాదములు}  ·  \te{అవును}  ·  \te{నమస్కారం}
\end{quote}

\begin{itemize}
\raggedright
  \item \emph{Trace it:} \te{వ} three times, saying \emph{va} as you finish each one
  \item \emph{Look:} back at any page of this chapter and find \te{వ} once more
\end{itemize}

\begin{tcolorbox}[breakable,colback=teal!4,colframe=teal!35!black,title={Before you move on}]
Which character is this — \te{వ}? What sound does it carry? (\emph{\textbf{va}}.) Name one word you already say that contains it.

One more, from much earlier: \te{౮} — what amount does it stand for?
\end{tcolorbox}
